\documentclass[11pt]{amsart}
\usepackage[a4paper,margin=1in]{geometry}
\usepackage{amsmath,amssymb,amsthm,mathtools}
\usepackage{booktabs,array,longtable,tabularx}
\usepackage{microtype}
\usepackage{enumitem}
\usepackage{etoolbox}
\usepackage[hidelinks]{hyperref}
\usepackage[nameinlink,noabbrev]{cleveref}
\usepackage[backend=biber,style=numeric,sorting=nyt,maxbibnames=99]{biblatex}
\addbibresource{references.bib}
\input{author_info.tex}

\newtheorem{theorem}{Theorem}[section]
\newtheorem{proposition}[theorem]{Proposition}
\newtheorem{lemma}[theorem]{Lemma}
\newtheorem{corollary}[theorem]{Corollary}
\theoremstyle{definition}
\newtheorem{definition}[theorem]{Definition}
\newtheorem{conjecture}[theorem]{Conjecture}
\theoremstyle{remark}
\newtheorem{remark}[theorem]{Remark}

\crefname{theorem}{theorem}{theorems}
\Crefname{theorem}{Theorem}{Theorems}
\crefname{proposition}{proposition}{propositions}
\Crefname{proposition}{Proposition}{Propositions}
\crefname{lemma}{lemma}{lemmas}
\Crefname{lemma}{Lemma}{Lemmas}
\crefname{corollary}{corollary}{corollaries}
\Crefname{corollary}{Corollary}{Corollaries}
\crefname{conjecture}{conjecture}{conjectures}
\Crefname{conjecture}{Conjecture}{Conjectures}

\newcommand{\Z}{\mathbb Z}
\newcommand{\Rset}{\mathcal R}
\newcommand{\OO}{\mathrm O}
\newcommand{\EE}{\mathrm E}

\title[Read-Once Factorial Divisibility Through Six Inputs]{The Read-Once Factorial-Divisibility Conjecture:\\Algebraic Proofs Through Six Inputs}
\author{\AuthorName}
\address{\AuthorAffiliation}
\email{\AuthorEmail}
\ifdefempty{\AuthorORCID}{}{\thanks{ORCID: \AuthorORCID}}
\thanks{Preprint DOI: \href{https://doi.org/10.5281/zenodo.21853500}{10.5281/zenodo.21853500}}
\date{August 2026}
\subjclass[2020]{11A07, 11B75, 05A05}
\keywords{factorial divisibility, read-once arithmetic expression, congruences, signed residue classes, divisibility frontier, conjecture}
\hypersetup{
  pdftitle={The Read-Once Factorial-Divisibility Conjecture: Algebraic Proofs Through Six Inputs},
  pdfauthor={\AuthorName},
  pdfsubject={Read-once arithmetic expressions and factorial divisibility},
  pdfkeywords={factorial divisibility, read-once expression, congruences, conjecture},
  pdfinfo={DOI={10.5281/zenodo.21853500}}
}

\begin{document}

\begin{abstract}
Let $S$ be a set of $n$ distinct integers.  The read-once
factorial-divisibility conjecture asks whether one can combine the elements of
$S$, using each element exactly once and using only addition, subtraction, and
multiplication, to obtain an element of $n!\Z$.  The conjecture remains open in
general.  We prove it algebraically for every $1\le n\le 6$.  The proofs use
congruences, signed residue classes, pigeonhole principles, and products of
expressions on disjoint blocks; exhaustive enumeration is not used as proof.
The crucial new ingredients for six inputs are a mixed-parity theorem stating
that three odd inputs and one even input always yield a multiple of $60$, and
an auxiliary theorem producing a $120$-multiple from four odd inputs coprime
to $15$.  These close the parity patterns $3\OO3\EE$ and $4\OO2\EE$.  We also
record current proved divisibility frontiers $F_n,O_n,E_n$ for arbitrary,
all-odd, and all-even inputs through $n=6$, emphasize that incomparable
moduli may require different expressions, and document the provenance and
dominance reductions behind those frontiers.
\end{abstract}

\maketitle
\pagestyle{plain}

\section{Introduction}
Arithmetic-expression construction usually asks whether prescribed numbers
can be combined to reach a particular target.  Here the target is a
congruence class: from $n$ distinct integers, every input must occur exactly
once as a leaf of one binary expression tree, and the root value should lie
in $n!\Z$.  The read-once requirement prevents reuse in identities such as
$(a-b)(a+b)$ and makes local congruence choices compete for leaves.  This
places the problem near arithmetic-expression construction
\cite{AlcockEtAl2020} and arithmetic read-once formulas
\cite{BshoutyHancockHellerstein1995}, while the repeated use of signed residue
classes is reminiscent of zero-sum methods
\cite{ErdosGinzburgZiv1961,GaoGeroldinger2006}.  The factorial target also
connects naturally with classical questions about factorials and generalized
factorials \cite{Bhargava2000}.

The central object of this paper is a conjecture, not a theorem for all
inputs: the statement for general $n$ remains open.  The contribution here is
to prove the conjecture by hand for the complete range $1\le n\le6$ and to
organize the small cases into reusable divisibility profiles.  A block need
not have one preferred modulus.  For example, three even integers
universally yield both an $18$-multiple and a $48$-multiple, generally by
different expressions.  Neither modulus divides the other, so either can be
the useful block in a larger construction.

\begin{theorem}[Main theorem]\label{thm:main}
For every integer $n$ with $1\le n\le6$ and every set of $n$ distinct
integers, there exists a read-once expression using only $+,-,\times$ whose
value belongs to $n!\Z$.
\end{theorem}

The cases through five inputs are developed from small congruence theorems.
The six-input case is divided by parity.  Five of the seven parity patterns
follow from earlier blocks.  The remaining patterns are closed by
\[
  3\OO1\EE\to 60
  \qquad\text{and}\qquad
  4\OO\ \text{with all inputs coprime to }15\to120.
\]
The main text keeps the structural proof.  Complete congruence branches are
placed in the appendices so that no case is omitted while the main argument
remains readable.

\section{Definitions and the conjecture}

\begin{definition}[Read-once expression]
Let $S=\{a_1,\dots,a_n\}\subset\Z$.  A read-once arithmetic expression on
$S$ is a fully parenthesized expression formed from the binary operations
$+,-,$ and $\times$, in which each input occurs exactly once.  Equivalently,
it is a full binary tree whose leaves are in bijection with $S$ and whose
internal nodes are labelled by $+,-,$ or $\times$.  Parentheses determine
the order of operations and do not count as additional operations.
\end{definition}

For example,
\[
 (a_1-a_2)(a_3+a_4)
\]
is a legal read-once expression on four inputs.  By contrast,
$a_1(a_2-a_3)+a_1$ is not legal because $a_1$ occurs twice, and
$(a_1-a_2)^2$ is not legal because exponentiation is not an allowed
operation and its expansion would reuse the inputs.

The value $0$ is allowed, because $0\in m\Z$ for every positive integer $m$.
The formal statements concern sets of distinct integers; repeated-input
cancellation is therefore not part of the problem.

\begin{conjecture}[Read-once factorial-divisibility conjecture]\label{conj:factorial}
For every integer $n\ge1$ and every set $S$ of $n$ distinct integers, there
exists a read-once expression on $S$, using only $+,-,\times$, whose value
lies in $n!\Z$.
\end{conjecture}

\begin{definition}[Guarantee notation]
Write $P_r(n;m)$ when every set of $n$ distinct integers having exactly $r$
odd elements admits a read-once expression divisible by $m$.  Write
$P(n;m)$ when the parity restriction is absent.
\end{definition}

For compact parity notation, $3\OO1\EE\to60$ means $P_3(4;60)$.  The
read-once expression may depend on the actual input set; the conjecture does
not require one fixed expression template to work for every set.

\begin{definition}[Recorded source-pool divisibility frontiers]
For fixed $n$, let $F_n$ denote the moduli maximal under divisibility among
the universal guarantees explicitly included in the source pool described in
\Cref{app:frontiers}; let $O_n$ and $E_n$ be the corresponding sets for the
all-odd and all-even source pools.  A modulus $m$ is removed when another
modulus $M$ in the same source pool satisfies $m\mid M$.  In that situation
we say that $M$ \emph{dominates} $m$, because the same expression that is a
multiple of $M$ is automatically a multiple of $m$.
\end{definition}

For example, once $63$ is proved for a class, the smaller guarantee $21$ is
not retained in its frontier because $21\mid63$.  This dominance deletion is
fundamentally different from combining guarantees obtained by different
expressions: knowing separately that $6$ and $8$ are guaranteed does not imply
that one expression is divisible by $\operatorname{lcm}(6,8)=24$.  The
recorded frontiers below summarize the displayed source pools; they are not
optimality claims and are not asserted to be closed under every further case
split that can be deduced from the proved parity theorems.

\section{General construction principles}

\begin{lemma}[Block product]\label{lem:block}
Suppose disjoint input blocks $A$ and $B$ admit read-once expressions
$E_A,E_B$ with $m\mid E_A$ and $q\mid E_B$.  Then $E_AE_B$ is a read-once
expression on $A\cup B$ divisible by $mq$.
\end{lemma}
\begin{proof}
The two subexpressions use disjoint leaves, and their product uses every leaf
in $A\cup B$ exactly once.
\end{proof}

\begin{lemma}[Two-input factor $3$]\label{lem:pair3}
Any two integers can be combined into a multiple of $3$.
\end{lemma}
\begin{proof}
If one input is divisible by $3$, use the product.  Otherwise both are
congruent to $\pm1\pmod3$, so their sum or difference is divisible by $3$.
\end{proof}

\begin{lemma}[Signed pairing modulo $2n$]\label{lem:signedpair}
Among any $n$ integers, either one input is divisible by $n$, or two inputs
$a,b$ satisfy $2n\mid a+b$ or $2n\mid a-b$.
\end{lemma}
\begin{proof}
Modulo $2n$, a residue $0$ or $n$ represents an integer divisible by $n$.
Otherwise the residues lie in the $n-1$ signed classes
$\{\pm1\},\dots,\{\pm(n-1)\}$.  Two of the $n$ inputs occupy one class, and
their suitable sum or difference is divisible by $2n$.
\end{proof}

\begin{lemma}[Sign normalization and transport]\label{lem:signnorm}
In a congruence argument one may replace each input $a_i$ by an independently
chosen signed representative $\varepsilon_i a_i$, where
$\varepsilon_i\in\{\pm1\}$.  More generally, let $E$ be any read-once
expression formed from the signed leaves $\varepsilon_i a_i$ using only the
binary operations $+,-,$ and $\times$.  Then there is a legal read-once
expression $E_0$ on the original leaves $a_i$, with the same underlying tree
and using only $+,-,$ and $\times$, such that $E=E_0$ or $E=-E_0$.  Hence
restoring the selected leaf signs does not change any divisibility conclusion.
\end{lemma}
\begin{proof}
Proceed by induction on the expression tree.  For a leaf $\varepsilon a$,
take the unsigned leaf $a$ and keep $\varepsilon$ only as an external sign
marker.  At a multiplication node, multiply the two external sign markers
and multiply the two unsigned child expressions.  At an addition or
subtraction node, factor out the external sign marker of the left child; the
relative sign of the right child determines whether that node is labelled
$+$ or $-$.  Thus every internal node remains one of the allowed binary
operations, every original leaf is used exactly once, and only one overall
sign remains at the root.  An overall sign does not affect divisibility.
\end{proof}

\begin{lemma}[Signed-class triangle]\label{lem:triangle}
Suppose $u,v,w$ lie in one signed residue class $\{\pm r\}$ modulo $m$,
where $r\not\equiv0\pmod m$ and $2r\not\equiv0\pmod m$.  For each pair,
let $\sigma_m(u,v)\in\{\pm1\}$ denote the unique sign satisfying
$m\mid u+\sigma_m(u,v)v$.  Then
\[
 \sigma_m(u,v)\sigma_m(v,w)\sigma_m(w,u)=-1.
\]
\end{lemma}
\begin{proof}
Write $u\equiv\varepsilon_ur$, $v\equiv\varepsilon_vr$ and
$w\equiv\varepsilon_wr\pmod m$, with each $\varepsilon_i\in\{\pm1\}$.
The condition $2r\not\equiv0\pmod m$ makes the admissible sign unique, and
$\sigma_m(u,v)=-\varepsilon_u\varepsilon_v$.  Multiplying the three edge
signs gives $-1$.
\end{proof}

The triangle relation will be used for $m=3,4,5$.  For $m=3$, all nonzero
residues form $\{\pm1\}$; for $m=4$, all odd residues form $\{\pm1\}$; and
for $m=5$, the nonzero residues form $\{\pm1\}$ and $\{\pm2\}$.

\section{Core blocks through five inputs}

This section records the blocks needed in the six-input proof.  Complete
small-case derivations, including complementary moduli used to form the
frontiers, appear in \Cref{app:small}.

\begin{theorem}[Three-input $12$ theorem]\label{thm:three12}
Every three integers admit a read-once expression divisible by $12$; that is,
$P(3;12)$.
\end{theorem}
\begin{proof}
See \Cref{app:three12}.  The proof separates the four parity patterns and
uses only congruences modulo $3$ and $4$.
\end{proof}

\begin{theorem}[Same-parity $18$ and even-input $48$]\label{thm:threeprofiles}
Any three integers of the same parity admit an expression divisible by $18$,
and any three even integers admit an expression divisible by $48$:
\[
 P_0(3;18),\quad P_3(3;18),\quad P_0(3;48).
\]
\end{theorem}
\begin{proof}
See \Cref{app:three18,app:three48}.
\end{proof}

\begin{theorem}[Four-input universal blocks]\label{thm:fourblocks}
Every four integers admit expressions divisible by $24$ and by $36$:
\[
 P(4;24),\qquad P(4;36).
\]
\end{theorem}
\begin{proof}
For $24$, apply \Cref{lem:signedpair} with $n=4$.  If one input is divisible
by $4$, multiply it by a $12$-multiple from the remaining three.  Otherwise
a pair yields an $8$-multiple, while the complementary pair yields a
$3$-multiple by \Cref{lem:pair3}.

For $36$, if one input is divisible by $3$, multiply it by a $12$-multiple
from the remaining three.  Assume henceforth that no input is divisible by
$3$.

If the number of odd inputs is $0$, $2$, or $4$, partition the four inputs
into two same-parity pairs.  In each pair choose a sum or difference divisible
by $3$.  Because the two entries of the pair have the same parity, that same
sum or difference is even, so each pair gives a $6$-block.  The two disjoint
blocks therefore give $36$.

If the parity pattern is $3\OO1\EE$, the three odd inputs give an
$18$-multiple by \Cref{thm:threeprofiles}, and multiplication by the remaining
even input supplies an additional factor $2$.  Hence the full expression is
divisible by $36$.

Finally consider $1\OO3\EE$.  Let $a$ be the odd input and $x,y,z$ the
even inputs.  Among $x,y,z$, two are congruent modulo $4$; relabel them
$x,y$.  Choose the sign for which $3\mid x\pm y$.  Since
$x\equiv y\pmod4$, both $x+y$ and $x-y$ are divisible by $4$, so the same
chosen block is divisible by $12$.  The remaining pair $a,z$ can be combined
into a multiple of $3$ by \Cref{lem:pair3}.  Multiplying the two disjoint
blocks gives a multiple of $36$.

These cases exhaust all parity patterns.
\end{proof}

\begin{theorem}[Four-input parity blocks]\label{thm:fourparity}
The following guarantees hold:
\[
\begin{aligned}
 &P_0(4;48),\ P_1(4;48),\ P_2(4;48),\ P_4(4;48),\ P_0(4;72),\\
 &P_1(4;40),\ P_2(4;40),\ P_3(4;60),\ P_4(4;60).
\end{aligned}
\]
\end{theorem}
\begin{proof}
See \Cref{app:four48,app:four72,app:oneodd40,app:twoodd40,app:fourodd60,app:mixed60}.
\end{proof}

\begin{theorem}[Five-input factorial target]\label{thm:five120}
Every five distinct integers admit a read-once expression divisible by
$120$; that is, $P(5;120)$.
\end{theorem}
\begin{proof}
Apply \Cref{lem:signedpair} with $n=5$.  If one input is divisible by $5$,
multiply it by a $24$-multiple from the remaining four.  Otherwise a pair
has a sum or difference divisible by $10$, and the remaining three inputs
yield a $12$-multiple.  Thus $10\cdot12=120$.
\end{proof}

\begin{theorem}[Five odd and five even inputs]\label{thm:fiveparity}
Five odd inputs always yield multiples of $180$ and $240$, and five even
inputs always yield a multiple of $480$:
\[
 P_5(5;180),\qquad P_5(5;240),\qquad P_0(5;480).
\]
\end{theorem}
\begin{proof}
See \Cref{app:fiveodd,app:fiveeven}.
\end{proof}

Consequently \Cref{conj:factorial} holds algebraically for $1\le n\le5$.
The six-input proof requires one additional four-leaf block.

\section{The mixed-parity modulus \texorpdfstring{$60$}{60}}

\begin{theorem}[Mixed-parity $60$ theorem]\label{thm:3o1e60}
Any three odd integers and one even integer admit a read-once expression
divisible by $60$:
\[
 3\OO1\EE\to60.
\]
\end{theorem}
\begin{proof}
See \Cref{app:mixed60} for the complete case analysis.
\end{proof}

\paragraph{Idea of proof.}
Let $a,b,c$ be odd and $e$ even.  The proof first disposes of an input
divisible by $5$, then separates the case $3\mid e$.  Under
$3\nmid e$, the three odd inputs contain either zero, one, or at least two
multiples of $3$, giving an exhaustive partition.  When no odd input is a
multiple of $3$, sign comparisons modulo $3$, $4$, and $5$, together with
the signed-class triangle relation, produce complementary blocks whose
moduli multiply to $60$.  When exactly one odd input is a multiple of $3$,
sign normalization modulo $5$ leaves exactly two possible signed-class
positions for $e$; the explicit modulo-$3$, modulo-$4$, and modulo-$5$
calculations are given in \Cref{app:caseE1,app:caseE2}.  Every simultaneous
divisibility claim there uses one and the same choice of sign.

\begin{theorem}[Four odd inputs coprime to $15$]\label{thm:4ocop120}
If four odd inputs are all coprime to $15$, then they admit a read-once
expression divisible by $120$.
\end{theorem}
\begin{proof}
Modulo $5$, the four inputs occupy two nonzero signed classes.  If the
class-size distribution is $3+1$ or $4+0$, take three inputs $u,v,w$ in one
class.  Were every edge to satisfy $\sigma_3=-\sigma_5$, the products of the
three edge signs modulo $3$ and modulo $5$ would be opposite.  By
\Cref{lem:triangle} both products equal $-1$, a contradiction.  Hence some
pair satisfies $\sigma_3=\sigma_5$, and its suitable sum or difference is
divisible by $30$.  The remaining odd pair yields a multiple of $4$.

If the distribution is $2+2$, write $a,b$ in one class and $c,d$ in the
other.  If either same-class pair, say $a,b$, satisfies
$\sigma_3(a,b)=\sigma_5(a,b)$, then the same signed sum or difference is
divisible by $2$, $3$, and $5$, so $a,b$ form a $30$-block; the remaining
odd pair $c,d$ forms a $4$-block.  Hence $30\cdot4=120$.

Assume therefore that both same-class pairs satisfy
$\sigma_3=-\sigma_5$.  For the pair $a,b$, the binary sign
$\sigma_4(a,b)$ must equal either $\sigma_3(a,b)$ or
$\sigma_5(a,b)$.  If $\sigma_4(a,b)=\sigma_3(a,b)$, then the same signed
sum or difference of $a,b$ is divisible by $12$, while the pair $c,d$, using
its modulo-$5$ sign, forms a $10$-block; thus $12\cdot10=120$.  If instead
$\sigma_4(a,b)=\sigma_5(a,b)$, then $a,b$ form a $20$-block, while $c,d$,
using their modulo-$3$ sign, form a $6$-block; thus $20\cdot6=120$.
\end{proof}

\section{Six inputs}

\begin{theorem}[Six-input factorial theorem]\label{thm:six720}
Every six distinct integers admit a read-once expression divisible by
$720=6!$.
\end{theorem}
\begin{proof}
We split by the number of odd inputs.  The complete leaf-by-leaf
constructions appear in \Cref{app:six}; the main exits are summarized below.

\begin{center}
\begin{tabular}{@{}cl@{}}
\toprule
Parity pattern & Divisibility construction \\
\midrule
$0\OO6\EE$ & $18\cdot10\cdot4=720$ \\
$1\OO5\EE$ & $20\cdot36$, or $10\cdot2\cdot36$ \\
$2\OO4\EE$ & $10\cdot72$, $20\cdot36$, $18\cdot40$, or equivalent branches \\
$3\OO3\EE$ & $6\cdot120$, or $12\cdot60$ \\
$4\OO2\EE$ & $6\cdot120$, $12\cdot60$, $20\cdot36$, or $120\cdot6$ \\
$5\OO1\EE$ & $60\cdot12$, or $20\cdot36$ \\
$6\OO0\EE$ & $18\cdot10\cdot4=720$ \\
\bottomrule
\end{tabular}
\end{center}

For $0\OO6\EE$ and $6\OO0\EE$, one first constructs an $18$-block,
then a $10$-block from two of the remaining four inputs, and finally a
$4$-block from the last pair.  The parity of each selected pair supplies the
factor $2$ accompanying the modulo-$9$ or modulo-$5$ divisibility.

For $1\OO5\EE$, an even multiple of $5$ supplies $10$; another even leaf
supplies $2$, and the remaining four leaves supply $36$.  If no even input
is divisible by $5$, two of the five even inputs collide in one of the four
nonzero even signed classes modulo $20$, producing a $20$-block, and the
remaining four leaves supply $36$.  The $5\OO1\EE$ argument is the odd
analogue: an odd multiple of $5$ is combined with a $12$-block to produce
$60$, while in the absence of such an input a signed-class collision modulo
$20$ produces $20$ and the complementary block produces $36$.

For $2\OO4\EE$, the proof first separates the presence of a multiple of
$5$.  Otherwise it checks whether the odd pair shares a signed class modulo
$5$, whether an even pair collides modulo $20$, and finally whether a
multiple of $3$ occurs.  If none does, a pair of even inputs gives $18$;
the unused leaves give either $10\cdot4$ or the previously proved
$2\OO2\EE\to40$.  These branches are written leaf by leaf in
\Cref{app:six2}.

For $3\OO3\EE$, if an even input is divisible by $3$, it supplies $6$ and
the remaining five inputs supply $120$.  Otherwise the three even inputs are
nonzero modulo $3$.  Two have the same residue modulo $4$; choose the sign
making their sum or difference divisible by $3$.  The same block is
divisible by $4$, hence by $12$.  The complementary $3\OO1\EE$ block
supplies $60$.

For $4\OO2\EE$, first dispose of an even multiple of $3$ by
$6\cdot120$.  If the two even inputs are congruent modulo $4$, they yield a
$12$-block and the four odd inputs yield $60$.  Otherwise write the even
inputs as $4\mid x$ and $y\equiv2\pmod4$.  An odd multiple of $3$ gives
$12\mid ax$ and leaves a $3\OO1\EE$ block of modulus $60$.  If no odd input
is divisible by $3$ but one is divisible by $5$, then $20\mid ax$ and the
remaining four inputs give $36$.  In the final case all four odd inputs are
coprime to $15$, so \Cref{thm:4ocop120} gives $120$; the even pair gives
$6$.

The other five parity patterns are proved in
\Cref{app:six0,app:six1,app:six2,app:six5,app:six6}.  The seven patterns
exhaust all possibilities.
\end{proof}

\begin{corollary}\label{cor:n6}
\Cref{conj:factorial} holds for every $1\le n\le6$.
\end{corollary}

\section{Recorded source-pool divisibility frontiers}

The following table records the nondominated moduli in the explicit source
pools used in this paper.  Each entry is a proved existence guarantee for
some expression; entries in the same row need not be attained simultaneously.
The table is a provenance summary rather than an exhaustive closure of every
parity-restricted consequence.

\small
\begin{longtable}{@{}c>{\raggedright\arraybackslash}p{0.285\textwidth}>{\raggedright\arraybackslash}p{0.285\textwidth}>{\raggedright\arraybackslash}p{0.285\textwidth}@{}}
\toprule
$n$ & $F_n$ (arbitrary) & $O_n$ (all odd) & $E_n$ (all even) \\
\midrule
\endfirsthead
\toprule
$n$ & $F_n$ (arbitrary) & $O_n$ (all odd) & $E_n$ (all even) \\
\midrule
\endhead
1 & $\{1\}$ & $\{1\}$ & $\{2\}$ \\
2 & $\{2,3\}$ & $\{3,4\}$ & $\{6,8\}$ \\
3 & $\{7,8,9,10,12\}$ & $\{8,10,12,14,18\}$ & $\{10,14,18,48\}$ \\
4 & $\{10,14,16,24,36\}$ & $\{14,36,48,60\}$ & $\{20,28,64,72,96\}$ \\
5 & $\begin{gathered}\{14,16,21,27,\\36,120\}\end{gathered}$ & $\begin{gathered}\{32,42,54,56,\\72,180,240\}\end{gathered}$ & $\begin{gathered}\{84,108,112,\\288,384,480\}\end{gathered}$ \\
6 & $\begin{gathered}\{49,56,63,64,70,\\81,84,96,100,108,720\}\end{gathered}$ & $\begin{gathered}\{100,112,140,168,192,\\196,216,252,324,720\}\end{gathered}$ & $\begin{gathered}\{100,140,196,252,324,\\512,672,720,864,960,2304\}\end{gathered}$ \\
\bottomrule
\end{longtable}
\normalsize

The products and dominance deletions producing these rows are listed in
\Cref{app:frontiers}.  Optimality is not claimed.

Beyond the source pool used in the factorial proof, we found further universal
four-input guarantees.  The additional constructions in
\Cref{app:four-refinements} prove $P(4;m)$ for
\[
 m\in\{17,19,21,22,23,25,26,27,28,30,32,33,35,40\}.
\]
Together with the previously recorded $24$- and $36$-guarantees, the
nondominated four-input universal moduli among the constructions displayed in
this manuscript are
\[
 \{17,19,21,22,23,24,25,26,27,28,30,32,33,35,36,40\}.
\]
The additional four-input refinements are not propagated into the five- and
six-input source-pool rows above.  A parity-by-parity guide to the factorial
constructions through six inputs is collected in \Cref{app:parity-summary}.

\section{Guaranteed moduli through six inputs}

To record one convenient universal modulus at each input size proved in this
paper, we may take
\[
 L_1=1,\quad L_2=2,\quad L_3=12,\quad L_4=24,\quad
 L_5=120,\quad L_6=720.
\]
Each $L_n$ is a separately proved guarantee; no optimality is asserted, and
no recurrence beyond six inputs is pursued here.  In particular, the general
factorial-divisibility conjecture remains open beyond the range covered by
\Cref{thm:main}.

\section{Exact residue dynamic programming as a checking tool}

For a fixed modulus $M$ and fixed inputs, exact read-once reachability can be
computed over residues.  For every nonempty subset $T$ of the input indices,
let $\Rset(T)$ be the residues modulo $M$ attained by expressions using
exactly the leaves in $T$.  Initialize singleton sets by the input residues.
For every bipartition $T=A\sqcup B$, combine $x\in\Rset(A)$ and
$y\in\Rset(B)$ by
\[
 x+y,\quad x-y,\quad y-x,\quad xy\pmod M.
\]
Induction on $|T|$ shows that this recurrence is complete.

The algorithm is not used as proof here.  It serves only as an independent
consistency check of delicate blocks.  The algebraic appendices, not finite
search, establish the stated theorems.

\section{Conclusion}
The read-once factorial-divisibility statement remains a conjecture for
general $n$.  This paper proves the complete range through six inputs.  The
decisive step is not exhaustive residue enumeration, but reusable congruence
blocks with complementary prime-power profiles, especially
$3\OO1\EE\to60$ and the four-odd coprime-to-$15$ modulus $120$.  Together
with the earlier three-, four-, and five-input results, these close all seven
parity patterns for $n=6$.  The first unresolved finite case is $n=7$.

\section*{Generative AI disclosure}
OpenAI ChatGPT (GPT-5.6) was used for exploratory mathematical discussion,
organization and revision of proof arguments, preparation and review of
auxiliary consistency-checking code, and English-language editing.  The tool
assisted in examining possible proof structures, identifying arguments
requiring clarification, and improving the presentation of the manuscript.
It was not used as a substitute for mathematical proof.  The author reviewed
all definitions, theorem statements, proofs, computations, references, and
conclusions and assumes full responsibility for the originality, accuracy,
and integrity of the manuscript.

\section*{Declarations}
\textbf{Funding.} No external funding was received for this work.

\textbf{Competing interests.} The author declares no competing interests.

\textbf{Data and code availability.} No external dataset is associated
with this work.  The LaTeX source accompanies the preprint record.  The
algebraic proofs are self-contained.  A supplementary Python script supplied
with the source independently reconstructs the residual signed-class lists
used in the $17/19$, $23$, and $35$ arguments and checks their displayed
congruence certificates; it serves only as a consistency check and is not
required for the proofs.

\appendix
\input{appendix_small_cases.tex}
\input{appendix_mixed60.tex}
\input{appendix_six_inputs.tex}
\input{appendix_parity_summary.tex}
\input{appendix_frontiers.tex}
\input{appendix_four_input_refinements.tex}

\printbibliography
\end{document}
